% SOURCE_AUTHORITY: sga5_fr_workpass.tex, lines 3283--3463 (LF-normalized unit).
% SOURCE_SHA256: 791F4EFFC5E02832D5D77ED03518C8156D6F07E4C8238B03545DB93D883FBB28
\begin{statement}{Corolario 4.5}
Bajo las hipótesis de 4.4, se tiene un cuadrado conmutativo
\[
\begin{tikzcd}[column sep=large,row sep=large]
\Hom(c_1'{}^*L_1,c_2'{}^!L_2)\otimes
\Hom(c_2''{}^*L_2,c_1''{}^!L_1)
 \arrow[r,"4.2.5"] \arrow[d,"3.7.6\otimes 3.7.6"'] &
\H^0(C,K_C)\arrow[d,"g_*"]\\
\Hom(d_1'{}^*f_{1*}L_1,d_2'{}^!f_{2*}L_2)\otimes
\Hom(d_2''{}^*f_{2*}L_2,d_1''{}^!f_{1*}L_1)
 \arrow[r,"4.2.5"'] &
\H^0(D,K_D),
\end{tikzcd}
\]
donde $g_*$ está definido por la flecha de adjunción
\[
g_*K_C=g_*g^!K_D\longrightarrow K_D.
\]
En otros términos, para
\[
u\in\Hom(c_1'{}^*L_1,c_2'{}^!L_2),\qquad
v\in\Hom(c_2''{}^*L_2,c_1''{}^!L_1),
\]
se tiene
\begin{equation}
\tag{4.5.1}
\langle f_*u,f_*v\rangle=g_*\langle u,v\rangle.
\end{equation}
\end{statement}

Para todo $S$-esquema propio $t:T\longrightarrow S$, denotemos por
\begin{equation}
\tag{4.6}
\int_T:\H^0(T,K_T)\longrightarrow \H^0(S,K_S)=\H^0(S,A)
\end{equation}
la flecha definida por la flecha de adjunción
\[
t_!K_T\longrightarrow A.
\]

\begin{statement}{Corolario 4.7 -- Fórmula de Lefschetz--Verdier}
Bajo las hipótesis de 4.4, supongamos que $Y_1=Y_2=S=D'=D''$. Para
\[
u\in\Hom(c_1'{}^*L_1,c_2'{}^!L_2),\qquad
v\in\Hom(c_2''{}^*L_2,c_1''{}^!L_1),
\]
se tiene
\begin{equation}
\tag{4.7.1}
\langle u_*,v_*\rangle=\int_C\langle u,v\rangle,
\end{equation}
donde
\[
u_*\in\Hom(f_{1*}L_1,f_{2*}L_2),\qquad
v_*\in\Hom(f_{2*}L_2,f_{1*}L_1)
\]
son los homomorfismos deducidos de $u$, $v$ por imagen directa sobre $S$ (3.7.6), y
\[
\langle u_*,v_*\rangle=\Tr(u_*v_*)=\Tr(v_*u_*)
\]
es el producto cup, en el sentido de (SGA 6 I 8.3), de los homomorfismos de complejos perfectos $u_*$, $v_*$ (4.1.5).
\end{statement}

Teniendo en cuenta 3.6.5, se tiene en particular:

\begin{statement}{Corolario 4.8}
Sean $t:T\longrightarrow S$ un $S$-esquema propio, $\Delta\subset T\times_S T$ la diagonal, $F:T\longrightarrow T$ un $S$-endomorfismo de $T$, $T^F=(\text{grafo de }F)\times_{(T\times_S T)}\Delta$ el esquema de puntos fijos de $F$, $L\in\ob\Dctf(T)$, $u\in\Hom(F^*L,L)$ una correspondencia de $L$ a $L$ con soporte en el grafo de $F$, $1\in\Hom(L,L)$ la correspondencia identidad, con soporte en $\Delta$ (3.2 b)). Entonces se tiene
\[
\Tr(u_*)=\int_{T^F}\langle u,1\rangle,
\]
donde $u_*$ es el endomorfismo de $t_*L$ compuesto de $t_*L\longrightarrow t_*F^*L$ (imagen inversa) y de $t_*u:t_*F^*L\longrightarrow t_*L$.
\end{statement}

De las observaciones de 3.2 a) y 3.6 se deduce, por otra parte:

\begin{statement}{Corolario 4.9}
Bajo las hipótesis de 4.7, supongamos que $X_i$ es liso sobre $S$, de dimensión relativa pura $r_i$, y que $C'$ (resp. $C''$) es un subesquema cerrado de $X$, finito y de dimensión Tor finita sobre $X_2$ (resp. $X_1$). Sean
\[
\cl(C')\in \H_{C'}^{2r_1}(X,A(r_1)),\qquad
\cl(C'')\in \H_{C''}^{2r_2}(X,A(r_2))
\]
las clases de cohomología de $C'$ y $C''$,
\[
\cl(C')_*\in\Hom(\R f_{1*}A,\R f_{2*}A),\qquad
\cl(C'')_*\in\Hom(\R f_{2*}A,\R f_{1*}A)
\]
los homomorfismos correspondientes (3.2 a), 3.7.6). Se tiene entonces
\[
\langle \cl(C')_*,\cl(C'')_*\rangle
=
\int_{C'\cap C''}\cl(C)\cl(D),
\]
donde
\[
\cl(C)\cl(D)\in
\H_{C'\cap C''}^{2(r_1+r_2)}(X,A(r_1+r_2))
=
\H^0(C'\cap C'',K_{C'\cap C''})
\]
es el producto cup habitual.
\end{statement}

Teniendo en cuenta 4.3, se tiene en particular:

\begin{statement}{Corolario 4.10}
Bajo las hipótesis de 4.9, supongamos que $C'\cap C''$ sea finito. Para $z\in C'\cap C''$, denotemos por $(C'\!\cdot C'')_z$ la imagen en $A$ de la multiplicidad de intersección de $C'$ y $C''$ en $z$. Entonces se tiene
\[
\langle \cl(C')_*,\cl(C'')_*\rangle
=
\sum_{z\in C'\cap C''}(C'\!\cdot C'')_z.
\]
\end{statement}

\begin{statement}{4.11}
Bajo las hipótesis de 4.9, se tiene en general
\begin{equation}
\tag{4.11.1}
\langle \cl(C')_*,\cl(C'')_*\rangle
=
\sum_{Z\in\pi_0(C'\cap C'')}(\cl(C')\cl(C''))_Z,
\end{equation}
donde $\pi_0(C'\cap C'')$ denota el conjunto de las componentes conexas de $C'\cap C''$, y $(\cl(C')\cl(C''))_Z$ la restricción de $\cl(C')\cl(C'')$ a $\H^0(Z,K_Z)$. Cuando $Z$ es una componente de dimensión $\geq 1$, el cálculo de $(\cl(C')\cl(C''))_Z$ plantea problemas, debido a que nos hallamos en una situación ``de exceso''. Si $Z$ es liso, de dimensión pura $d$, si $C'$ y $C''$ son lisos en un entorno de $Z$, y si $X$ es proyectivo, puede demostrarse que se tiene
\begin{equation}
\tag{4.11.2}
\int_Z(\cl(C')\cl(C''))_Z=\int_Z c_d(N_Z),
\end{equation}
donde $N_Z$ es el haz localmente libre de rango $d$ definido por la sucesión exacta de fibrados normales
\[
0\longrightarrow N_{Z/C'}\oplus N_{Z/C''}\longrightarrow N_{Z/X}\longrightarrow N_Z\longrightarrow 0,
\]
y $c_d(N_Z)\in \H^{2d}(Z,A(d))$ denota su $d$-ésima clase de Chern (SGA 5 VII) (utilizar (SGA 6 XIV 6.3) para reducirse a un cálculo de un número de intersección en teoría K, y aplicar (SGA 6 VII 2.5), cf. ([10] 4.1); la hipótesis de proyectividad es probablemente innecesaria).

Sea $T$ un $S$-esquema proyectivo liso; tomemos como correspondencia $C'$ (resp. $C''$) el grafo de un $S$-endomorfismo $F$ (resp. la diagonal), y supongamos que el esquema de puntos fijos $T^F=C'\cap C''$ sea liso. El haz $N_{(T^F)}$ anterior no es sino el haz tangente a $T^F$, de modo que 4.11.2 implica, por la fórmula de Gauss--Bonnet (SGA 5 VII 4.9), el análogo de una fórmula bien conocida por los topólogos:
\begin{equation}
\tag{4.11.3}
\Tr(F,\R\Gamma(T,A))=\EP(T^F),
\end{equation}
donde $\EP(T^F)$ denota la característica de Euler--Poincaré de $T_s^F$ para $s$ un punto geométrico sobre $s$ (loc. cit.).
\end{statement}

El cálculo de $(\cl(C')\cl(C''))_Z$ para una componente $Z$ de dimensión $\geq 1$, no lisa, no se ha abordado en general, al menos hasta donde sabe el redactor.

\begin{statement}{4.12}
En la situación de 4.7, sea $z$ un punto aislado de $C$, con proyección $z_i$ sobre $X_i$, $a'$ (resp. $a''$) sobre $C'$ (resp. $C''$), y sea $\bar z$ un punto geométrico sobre $z$, de imagen $\bar z_i$ sobre $z_i$. Supongamos que $c_2'$ (resp. $c_1''$) es cuasifinito y plano en $a'$ (resp. $a''$), que $L_i$ es perfecto en un entorno de $z_i$, y que $u$ (resp. $v$) viene dado en un entorno de $a'$ (resp. $a''$) por
\[
\Tr_{c_2'}:c_{2!}'c_2'{}^*L_2\longrightarrow L_2
\quad\text{y}\quad
s\in\Hom(c_1'{}^*L_1,c_2'{}^*L_2)
\]
(resp.
\[
\Tr_{c_1''}:c_{1!}''c_1''{}^*L_1\longrightarrow L_1
\quad\text{y}\quad
t\in\Hom(c_2''{}^*L_2,c_1''{}^*L_1)).
\]
Denotemos por $s_{\bar z}\in\Hom(L_{1\bar z_1},L_{2\bar z_2})$ (resp. $t_{\bar z}\in\Hom(L_{2\bar z_2},L_{1\bar z_1})$) el homomorfismo deducido de $s$ (resp. $t$) por paso a las fibras. Denotemos además, como en 3.2 a), por $\cl(c')$ (resp. $\cl(c'')$) la correspondencia de $X_1$ a $X_2$ (resp. de $X_2$ a $X_1$) definida, en un entorno de $a'$ (resp. $a''$), por $c_2'$ (resp. $c_1''$), y por $(C'\!\cdot C'')_z$ el término local $\langle\cl(c'),\cl(c'')\rangle_z$ correspondiente (igual a la imagen en $A$ de la multiplicidad de intersección de $C'$ y $C''$ en $z$ cuando $C'$ y $C''$ son subesquemas cerrados de $X$, con $X_1$ y $X_2$ lisos). Se tiene entonces:
\begin{equation}
\tag{4.12.1}
\langle u,v\rangle_z=(C'\!\cdot C'')_z\langle s_{\bar z},t_{\bar z}\rangle.
\end{equation}
Esto resulta, en efecto, fácilmente de 4.2.6.
\end{statement}

\begin{statement}{4.13}
Si $T$ es un esquema sobre $S$, denotemos por $DF(T)$ la categoría derivada filtrada de los complejos de $A_T$-módulos con filtración finita ([9] V 1), y por $DF_{\ctf}(T)$ la subcategoría plena formada por los $L$ tales que $\gr(L)\in\ob\Dctf(T)$. El sorites de compatibilidad de los funtores derivados filtrados con el paso a los graduados asociados implica que $DF_{\ctf}$ es estable por las seis operaciones $f^*$, $f_*$, $f_!$, $f^!$, $\otimes^{\mathbf L}$, $\uRHom$. Las definiciones y los resultados de los números anteriores son válidos, mutatis mutandis, en $DF_{\ctf}$. El apareamiento 4.1.2 es compatible con el paso a los graduados asociados (considerándose $K_X$ provisto de la filtración trivial), y de ello resulta que, con las notaciones de 4.2.5, para $L_i\in\ob DF_{\ctf}(X_i)$,
\[
u\in\Hom_{DF}(c_1^*L_1,c_2^!L_2),\qquad
v\in\Hom_{DF}(d_2^*L_2,d_1^!L_1),
\]
se tiene la siguiente igualdad en $\H^0(E,K_E)$:
\begin{equation}
\tag{4.13.1}
\langle u,v\rangle=\sum_i\langle \gr^i u,\gr^i v\rangle.
\end{equation}
Como observa Deligne, esta observación se aplica en particular a la situación considerada por Langlands en ([13] §7). Las correspondencias que considera en (loc. cit. 7.11) preservan, en efecto, cerca de los puntos fijos, filtraciones finitas de los haces, cuyos graduados asociados inducen, sobre la localización estricta en la imagen de un punto fijo, bien un haz constante, bien la prolongación por cero de un haz constante sobre el complemento del punto cerrado. El cálculo de los términos locales se reduce fácilmente, por 4.13, a 4.12.1 y al caso de haces concentrados en los puntos cerrados que son imágenes de los puntos fijos, caso que se trata directamente. Aplicando 4.7, se obtiene la fórmula de Lefschetz requerida en ([13] 7.11). Señalemos que el enunciado general (loc. cit. 7.12) es cierto; la demostración, más delicada, se dará en III B.
\end{statement}



\section*{5. Complementos}
